\documentclass[10pt,letterpaper]{article}
\usepackage{amsmath,amssymb,amsthm}
\begin{document}
\pagestyle{empty}
\section{Modular arithmetic}
Two integers are congruent modulo $m$ if $m$ divides their difference. The residues form a ring $\mathbb{Z}/m\mathbb{Z}$.

\begin{tabular}[t]{|l|p{5cm}|}
\hline
Operation & Rule\\ \hline
Sum & $a+b \bmod m$ is well defined on classes.\\[4pt]
Product & Likewise for $ab$; this is why we can reduce at each step.\\ \hline
\multicolumn{2}{|c|}{Division requires $\gcd(a,m)=1$}\\ \hline
\end{tabular}

\bigskip
\begin{description}
\item[Chinese remainder theorem] If $\gcd(m,n)=1$ then $\mathbb{Z}/mn\cong\mathbb{Z}/m\times\mathbb{Z}/n$.
\item[Euler] $a^{\varphi(m)}\equiv1\pmod m$ when $\gcd(a,m)=1$.
\end{description}

\begin{quotation}
Congruences let us do arithmetic on a clock face, and the clock face is a finite world where every nonzero element of a prime clock has an inverse.
\end{quotation}

\raggedright
Solving $7x\equiv3\pmod{11}$: since $7\cdot8=56\equiv1$, we get $x\equiv24\equiv2$. \vrule width 0.4pt height 8pt\ Check: $7\cdot2=14\equiv3$.
\end{document}
